% บทคัดย่อภาษาไทย
% หมายเหตุ: ไม่ควรเกิน 300 คำ
% ควรครอบคลุม: วัตถุประสงค์ → วิธีการ → ผลลัพธ์ → การประยุกต์ใช้

โครงงานนี้มีวัตถุประสงค์เพื่อศึกษาและพัฒนาระบบจำแนกและวิเคราะห์ข่าวสารอัตโนมัติด้วยเทคนิคการประมวลผลภาษาธรรมชาติแบบเฉพาะประเภท (Category-Specific NLP Analytics) โดยระบบพัฒนาด้วยภาษา Python ร่วมกับเฟรมเวิร์ค FastAPI สำหรับส่วน Backend และภาษา JavaScript ร่วมกับ Tailwind CSS และ Socket.IO สำหรับส่วน Frontend ระบบทำการรวบรวมข่าวจากเว็บไซต์ข่าวไทยด้วยเทคนิคเว็บสแครปปิง การอ่าน RSS Feed และการจำลองเบราว์เซอร์ด้วย Playwright จากนั้นส่งเนื้อหาข่าวในรูปแบบ Markdown ไปยังโมเดลภาษาขนาดใหญ่ผ่านบริการ KKU Gen.AI เพื่อสร้างบทสรุป ประเด็นสำคัญ อารมณ์ของข่าว และคำสำคัญในภาษาเดียวกับต้นฉบับ พร้อมจำแนกหมวดหมู่ข่าวออกเป็นหมวดหมู่เฉพาะทางด้วยการตัดคำภาษาไทยด้วย PyThaiNLP ร่วมกับกฎคำสำคัญและโมเดล SVM และแสดงผลข่าวพร้อมสรุปบนเว็บแอปพลิเคชันแบบเรียลไทม์ผ่าน WebSocket ระบบยังมีกลไกการจัดการแคชและคุกกี้เพื่อปรับแต่งประสบการณ์ผู้ใช้และป้อนข้อมูลให้อัลกอริทึมแนะนำข่าว โครงงานนี้เป็นประโยชน์ต่อการลดภาระของผู้อ่านในการติดตามข่าวจากหลายเว็บไซต์ และเป็นแนวทางในการประยุกต์ใช้เทคนิค NLP และโมเดลภาษาขนาดใหญ่สำหรับการวิเคราะห์ข่าวสารอัตโนมัติ